\documentclass[11pt]{article}

\usepackage[T1]{fontenc}
\usepackage[utf8]{inputenc}
\usepackage[ngerman]{babel}
\usepackage{amsmath, amssymb}
\usepackage{geometry}
\usepackage[hidelinks]{hyperref}
\usepackage{xcolor}
\usepackage{enumitem}
\usepackage{listings}

\geometry{a4paper, margin=1in}

\lstdefinestyle{leanstyle}{%
  basicstyle=\ttfamily\small,
  breaklines=true,
  columns=fullflexible,
  frame=single,
  rulecolor=\color{black!25},
  backgroundcolor=\color{black!3},
  keywordstyle=\color{blue!60!black},
  commentstyle=\color{green!40!black},
  showstringspaces=false
}

\lstset{style=leanstyle}

\title{Anmerkungen zum Lean-Code\\
\large Cut-and-Project-Periodenformel}
\author{Review-Kommentare}
\date{\today}

\begin{document}
\maketitle

\section*{Gesamturteil}

Der Lean-Code ist beeindruckend und formalisiert einen erheblichen Teil der kombinatorischen Kernstruktur des Papers. Besonders stark sind die Formalisierung der Restklassenverteilung, der zyklischen Intervalle, der Minimalitätsargumente und der set-valued Dichotomie.

Gleichzeitig sollte der Code in seiner jetzigen Form nicht ohne Präzisierung als vollständige Formalisierung der geometrischen Cut-and-Project-Konstruktion präsentiert werden. Der wichtigste Punkt ist die Modelltreue: Der Code formalisiert im Wesentlichen ein normalisiertes Restklassenmodell, während die tatsächlichen geometrischen Restklassen im Paper durch
\[
    c_r \equiv -\alpha\beta^{-1}r \pmod D
\]
gegeben sind.

\section{Stärken des Codes}

Die folgenden Bestandteile sind mathematisch und formal sehr gelungen:

\begin{itemize}[leftmargin=2em]
    \item \texttt{coprime\_alpha\_D} und \texttt{coprime\_beta\_D} bilden die elementare Zahlentheorie sauber ab.
    \item \texttt{count\_hits} ist eine gute Formalisierung der Trefferanzahl in einem Block aufeinanderfolgender Restklassen.
    \item \texttt{count\_hits\_lt\_D}, \texttt{count\_hits\_ge\_D}, \texttt{uniform\_residue\_distribution} und \texttt{non\_uniform\_residue\_distribution} treffen genau die zentrale Fallunterscheidung.
    \item \texttt{cyclic\_interval\_stabilizer\_trivial} formalisiert den entscheidenden Minimalitätsschritt.
    \item \texttt{generic\_minimality} ist konzeptionell überzeugend: Aus einem Periodenkandidaten entsteht eine Restklassen-Translation; diese stabilisiert die schwere Menge; der triviale Stabilisator erzwingt die behauptete Minimalität.
    \item Die set-valued Konstruktion über \texttt{set\_indicator}, \texttt{set\_size}, \texttt{set\_sorted} und \texttt{set\_difference\_sequence} spiegelt die mathematische Idee klar wider.
\end{itemize}

Der Code ist damit keine bloße Skizze, sondern eine ernsthafte formale Begleitung des mathematischen Arguments.

\section{Wichtigster kritischer Punkt: Der Multiplikator wird nicht durchgezogen}

Im Paper sind die tatsächlichen projektiven Restklassen
\[
    c_r \equiv -\alpha\beta^{-1}r \pmod D.
\]
Im Lean-Code werden zwar

\begin{lstlisting}
def multiplier ...
def residue_bijection ...
\end{lstlisting}

definiert, aber der Hauptteil arbeitet später mit

\begin{lstlisting}
def count_hits (D : Nat) [NeZero D] (r0 N : Nat) (x : ZMod D) : Nat :=
  (Finset.range N).filter (fun i => (r0 + i : ZMod D) = x) |>.card
\end{lstlisting}

Das entspricht den Restklassen
\[
    r_0, r_0+1, \ldots, r_0+N-1 \pmod D,
\]
nicht unmittelbar den geometrischen Restklassen
\[
    m r_0, m(r_0+1), \ldots, m(r_0+N-1) \pmod D,
    \qquad m=-\alpha\beta^{-1}.
\]

Mathematisch ist das für die Periodenlänge wahrscheinlich unschädlich, weil Multiplikation mit einer Einheit modulo~\(D\) nur eine Bijektion der Restklassen ist. Aber die tatsächlich sortierte Gap-Folge kann durch diese Multiplikation anders aussehen. Daher sollte entweder der Multiplikator in der konkreten Enumeration formal eingebaut werden, oder die Lean-Sektion im Paper sollte klar sagen, dass ein normalisiertes kombinatorisches Restklassenmodell formalisiert wurde.

\section{Vorschlag: \texttt{count\_hits\_unit}}

Eine mögliche Verbesserung wäre eine allgemeinere Trefferfunktion:

\begin{lstlisting}
def count_hits_unit
    (D : Nat) [NeZero D]
    (u : (ZMod D)^*)
    (r0 N : Nat)
    (x : ZMod D) : Nat :=
  (Finset.range N).filter
    (fun i => u.val * (r0 + i : ZMod D) = x) |>.card
\end{lstlisting}

Für das geometrische Modell wäre dann

\begin{lstlisting}
u = multiplier alpha beta h
\end{lstlisting}

und die tatsächliche Restklassenabbildung wäre formal enthalten.

Dann sollte man ein Transferlemma beweisen:

\begin{lstlisting}
count_hits_unit D u r0 N x =
count_hits D r0 N (u^{-1}.val * x)
\end{lstlisting}

Damit könnten viele vorhandene Beweise weiterverwendet werden. Der schwere Satz wäre dann nicht mehr direkt in der natürlichen Restklassenordnung ein Intervall, sondern nach Rücktransport durch \(u^{-1}\). Das entspricht genau der Argumentation im Paper.

\section{Benennung und Reichweite der Formalisierung}

Der Name

\begin{lstlisting}
main_theorem_concrete
\end{lstlisting}

wirkt etwas stärker, als die aktuelle Modellierung rechtfertigt. Er klingt so, als sei die volle geometrische Konstruktion formalisiert. Treffender wären Namen wie:

\begin{lstlisting}
main_theorem_residue_model
main_theorem_normalized_residue_model
\end{lstlisting}

Ähnliches gilt für die Typeclass

\begin{lstlisting}
class GeometricProjection ...
\end{lstlisting}

Die Axiome sind eher periodisch-kombinatorisch als geometrisch. Ein präziserer Name wäre zum Beispiel:

\begin{lstlisting}
class ResidueProjectionModel ...
class PeriodicResidueModel ...
\end{lstlisting}

Wenn der Name \texttt{GeometricProjection} beibehalten wird, sollte der Kommentar deutlicher sein, etwa:

\begin{lstlisting}
/--
Abstract interface for the combinatorial residue model obtained after
reducing the geometric cut-and-project construction.
-/
\end{lstlisting}

\section{Umbenennung von \texttt{non\_uniform\_residue\_distribution}}

Das Theorem

\begin{lstlisting}
theorem non_uniform_residue_distribution ...
\end{lstlisting}

gilt offenbar auch im Fall \(s=0\), also im uniformen Fall \(D\mid N\). Denn dann ist die Anzahl der \(q+1\)-Klassen gleich \(0\), und die Anzahl der \(q\)-Klassen gleich \(D\).

Der Name ist daher etwas irreführend. Besser wäre:

\begin{lstlisting}
theorem residue_distribution
theorem quotient_remainder_residue_distribution
\end{lstlisting}

Daraus könnten dann zwei Spezialisierungen abgeleitet werden:

\begin{lstlisting}
theorem non_uniform_residue_distribution_of_not_dvd ...
theorem uniform_residue_distribution ...
\end{lstlisting}

\section{Der degenerierte Fall ist im Code stärker als im Paper}

Der Code beweist in

\begin{lstlisting}
lemma period_degenerate_concrete ...
\end{lstlisting}

tatsächlich die Minimalität der Periode \(N/D\). Das ist sehr gut. Im Paper sollte der analoge Beweisschritt ebenfalls explizit ergänzt werden: Die Folge
\[
    \underbrace{0,\ldots,0}_{N/D-1},1,
    \underbrace{0,\ldots,0}_{N/D-1},1,\ldots
\]
hat nicht nur Periode \(N/D\), sondern minimale Periode \(N/D\). Für \(N/D>1\) liegt die \(1\) genau einmal pro Block der Länge \(N/D\), sodass keine kleinere positive Periode möglich ist.

\section{Architektur von \texttt{main\_theorem}}

Das Theorem

\begin{lstlisting}
theorem main_theorem ...
  [GeometricProjection alpha beta omega (difference_sequence alpha beta omega)] :
\end{lstlisting}

ist eine Mischform: Es verwendet abstrakte Axiome, aber bereits die konkrete Sequenz \texttt{difference\_sequence}. Architektonisch sauberer wäre eine abstrakte Version:

\begin{lstlisting}
theorem main_theorem_abstract
    (seq : Z -> Z)
    [GeometricProjection alpha beta omega seq] :
    ...
    HasPeriodLength seq L
\end{lstlisting}

und anschließend ein konkreter Wrapper:

\begin{lstlisting}
theorem main_theorem_concrete :
    HasPeriodLength (difference_sequence alpha beta omega) L
\end{lstlisting}

Das ist kein mathematischer Fehler, aber es würde die Struktur des Codes deutlicher machen.

\section{Modulo durch \texttt{set\_size}}

Die Definition

\begin{lstlisting}
def set_sorted ... :=
  let M := (set_size alpha beta omega : Z)
  let r := (i % M).toNat
  ...
\end{lstlisting}

verwendet Modulo durch \texttt{set\_size}. Lean ist total und erlaubt diese Definition auch bei \(M=0\). In den relevanten Fällen wird später bewiesen, dass \texttt{set\_size = N} oder \texttt{set\_size = D} gilt und damit positiv ist. Mathematisch ist das also in den Theoremen kontrolliert.

Für bessere Lesbarkeit und Robustheit wäre aber ein Lemma sinnvoll:

\begin{lstlisting}
lemma set_size_pos ...
\end{lstlisting}

Dieses Lemma könnte explizit dokumentieren, dass die set-valued Enumeration in den relevanten Fällen wohldefiniert als Periodenblock der positiven Länge ist.

\section{Heartbeat-lastige Beweise}

Mehrere Beweise verwenden hohe Heartbeat-Grenzen, zum Beispiel:

\begin{lstlisting}
set_option maxHeartbeats 400000
set_option maxHeartbeats 800000
\end{lstlisting}

Das ist für ein Forschungsprojekt akzeptabel, aber für langfristige Wartbarkeit fragil. Besonders \texttt{sorted\_multiset\_add\_q} wirkt technisch schwer. Es wäre besser, die dort verwendeten Aussagen über Division, Modulo und den Fall \(N=qD\) in kleinere allgemeine Lemmas auszulagern.

\section{Set-valued Teil}

Der set-valued Teil ist sehr überzeugend strukturiert:

\begin{itemize}[leftmargin=2em]
    \item Phase A: strukturelle Dichotomie über \texttt{count\_hits\_lt\_D} und \texttt{count\_hits\_ge\_D}.
    \item Phase B: abstrakter Satz \texttt{set\_main\_theorem}.
    \item Phase C: konkrete Konstruktion über \texttt{set\_indicator}, \texttt{set\_cumulative\_hits}, \texttt{set\_V}, \texttt{set\_size}, \texttt{set\_sorted} und \texttt{set\_difference\_sequence}.
\end{itemize}

Das ist didaktisch klar und mathematisch passend: Multiplizitäten werden auf \(\{0,1\}\) reduziert.

\section{Formulierungsvorschlag für das Paper}

Mit dem Code in der jetzigen Form wäre folgende Beschreibung angemessen:

\begin{quote}
The Lean formalisation verifies the algebraic and combinatorial core of the argument in a normalised residue model. It formalises the coprimality lemmas, the uniform and non-uniform residue-distribution statements, the cyclic-interval stabiliser argument, the generic minimality proof, the degenerate case, and the set-valued dichotomy. The Euclidean cut-and-project geometry is not formalised from first principles; instead, the formalisation starts from the residue model obtained after reducing the geometry to arithmetic progressions modulo \(D\).
\end{quote}

Falls der Multiplikator \(m=-\alpha\beta^{-1}\) in die konkrete Enumeration eingebaut wird, kann diese Formulierung entsprechend stärker ausfallen.

\section*{Schlussurteil}

Der Lean-Code ist sehr gut als kombinatorische Formalisierung. Er enthält die entscheidenden Minimalitätsargumente und stützt das Paper substanziell.

Der zentrale Vorbehalt lautet:

\begin{quote}
Der Code formalisiert aktuell nicht sichtbar die tatsächliche geometrische Restklassenabbildung
\[
    r \mapsto -\alpha\beta^{-1}r \pmod D,
\]
sondern arbeitet im Wesentlichen mit dem normalisierten Fall \(m=1\).
\end{quote}

Empfehlung:

\begin{enumerate}[leftmargin=2em]
    \item Für arXiv kurzfristig: Die Lean-Sektion vorsichtig formulieren und von einem normalisierten Restklassenmodell sprechen.
    \item Für maximale Qualität: \texttt{count\_hits\_unit} einführen und den Multiplikator \(m\) tatsächlich durch die konkrete Enumeration ziehen.
    \item Für Codepflege: \texttt{non\_uniform\_residue\_distribution} umbenennen, \texttt{main\_theorem} abstrakter machen und heartbeat-lastige Beweise modularisieren.
\end{enumerate}

Dann wäre die Lean-Formalisation eine sehr solide und glaubwürdige Begleitung des Papers.

\end{document}
